\documentclass[12pt,a4paper]{article}
\usepackage{amsmath, amssymb}

\begin{document}
	
	\section*{Formal Definition of Transfer Function as a Ratio}
	
	We want to formalize the conditions under which the ratio
	$$
	G(s) = \frac{Y(s)}{U(s)}
	$$
	is well-defined and meaningful.
	
	\subsection*{1. The setting: LTI systems}
	We consider linear time-invariant (LTI) systems described by ordinary differential equations with constant coefficients:
	$$
	a_n \frac{d^n y(t)}{dt^n} + a_{n-1}\frac{d^{n-1}y(t)}{dt^{n-1}} + \dots + a_0 y(t) 
	= b_m \frac{d^m u(t)}{dt^m} + b_{m-1}\frac{d^{m-1}u(t)}{dt^{m-1}} + \dots + b_0 u(t),
	$$
	where $u(t)$ is the input, $y(t)$ is the output, and the coefficients $a_i, b_j \in \mathbb{R}$.
	
	This is the most general finite-dimensional continuous-time LTI system.
	
	\subsection*{2. Applying the Laplace transform}
	Taking the Laplace transform with \emph{zero initial conditions}:
	$$
	\mathcal{L}\!\left\{ \frac{d^k y(t)}{dt^k} \right\} = s^k Y(s), \quad 
	\text{if all } y(0^-), y'(0^-), \dots = 0.
	$$
	
	Substituting into the differential equation gives:
	$$
	(a_n s^n + a_{n-1} s^{n-1} + \dots + a_0) Y(s) 
	= (b_m s^m + b_{m-1} s^{m-1} + \dots + b_0) U(s).
	$$
	
	\subsection*{3. Defining the ratio}
	Rearranging:
	$$
	\frac{Y(s)}{U(s)} = 
	\frac{b_m s^m + b_{m-1} s^{m-1} + \dots + b_0}
	{a_n s^n + a_{n-1} s^{n-1} + \dots + a_0}.
	$$
	
	We define:
	$$
	G(s) = \frac{Y(s)}{U(s)}.
	$$
	
	Thus, $G(s)$ is the ratio of two polynomials in $s$, uniquely determined by the system’s coefficients.
	
	\subsection*{4. Conditions for validity of the ratio}
	The ratio $G(s)$ is meaningful only if the following conditions hold:
	
	\begin{enumerate}
		\item \textbf{Linearity:} the system satisfies superposition.
		\item \textbf{Time-invariance:} coefficients $a_i, b_j$ are constants.
		\item \textbf{Causality:} the system does not depend on future inputs.
		\item \textbf{Zero initial conditions:} otherwise extra terms appear in $Y(s)$.
		\item \textbf{Laplace existence:} input and output must be of exponential order.
		\item \textbf{Denominator nonzero:} $a_n \neq 0$, and for each $s$ considered, $a_n s^n + \dots + a_0 \neq 0$.
	\end{enumerate}
	
	\subsection*{5. Meaning of the ratio}
	With these conditions, the system reduces to
	$$
	Y(s) = G(s)\, U(s).
	$$
	
	So the ratio $G(s)$ is not merely ``output divided by input'' in a naive sense, but the \emph{operator} that converts any Laplace-transformed input into the corresponding output.
	
	\medskip
	
	In the time domain:
	$$
	y(t) = (g * u)(t),
	$$
	where $g(t)$ is the impulse response.  
	In the Laplace domain, convolution becomes multiplication, and $G(s)$ is exactly the multiplier.
	
	\subsection*{6. Formal statement}
	For a causal, linear time-invariant, finite-dimensional system with zero initial conditions and inputs/outputs of exponential order, the ratio
	$$
	G(s) = \frac{Y(s)}{U(s)}
	$$
	is well-defined, unique, and represents the system’s \textbf{transfer function}. It encodes the system’s dynamics completely (up to initial conditions).
	
	\bigskip
	
	\textbf{In words:} The transfer function is the precise mathematical rule that tells you, in the Laplace domain, how much of the input signal (in amplitude and phase) passes through the system to produce the output.
	
	\section*{Laplace Transform Emergency Kit}
	
	\subsection*{1. Definition}
	For a function $f(t)$, $t \geq 0$:
	$$
	\mathcal{L}\{f(t)\} = F(s) = \int_0^{\infty} f(t) e^{-st}\, dt,
	$$
	where $s = \sigma + j\omega$ is a complex variable.
	
	The exponential $e^{-st}$ is the kernel that converts a time-domain signal into its representation in the $s$-domain.  
	Existence requires $f(t)$ be of exponential order: 
	$$
	|f(t)| \leq M e^{\alpha t}, \quad \text{for some constants } M,\alpha.
	$$
	
	\subsection*{2. Purpose}
	\begin{itemize}
		\item In the time domain, systems are described by differential equations.
		\item In the Laplace domain, differentiation and integration become multiplication and division by $s$.
		\item This converts ordinary differential equations into algebraic equations, which are easier to solve.
		\item System behavior (stability, poles, resonance) is naturally represented in the $s$-domain.
	\end{itemize}
	
	\subsection*{3. Core operational properties}
	\begin{align*}
		\mathcal{L}\{af(t) + bg(t)\} &= aF(s) + bG(s), \\
		\mathcal{L}\{f'(t)\} &= sF(s) - f(0^+), \\
		\mathcal{L}\{f''(t)\} &= s^2 F(s) - sf(0^+) - f'(0^+), \\
		\mathcal{L}\left\{\int_0^t f(\tau)\, d\tau\right\} &= \frac{1}{s}F(s), \\
		\mathcal{L}\{(f*g)(t)\} &= F(s)G(s).
	\end{align*}
	
	\subsection*{4. Common transforms}
	
	\[
	\begin{array}{|c|c|c|}
		\hline
		f(t) & F(s) & \text{Region of Convergence} \\
		\hline
		1 & \tfrac{1}{s} & \Re(s) > 0 \\
		t & \tfrac{1}{s^2} & \Re(s) > 0 \\
		t^n & \tfrac{n!}{s^{n+1}} & \Re(s) > 0 \\
		e^{at} & \tfrac{1}{s-a} & \Re(s) > a \\
		\sin(\omega t) & \tfrac{\omega}{s^2+\omega^2} & \Re(s) > 0 \\
		\cos(\omega t) & \tfrac{s}{s^2+\omega^2} & \Re(s) > 0 \\
		e^{at}\sin(\omega t) & \tfrac{\omega}{(s-a)^2+\omega^2} & \Re(s) > a \\
		e^{at}\cos(\omega t) & \tfrac{s-a}{(s-a)^2+\omega^2} & \Re(s) > a \\
		\delta(t) & 1 & \text{all $s$} \\
		u(t) & \tfrac{1}{s} & \Re(s) > 0 \\
		\hline
	\end{array}
	\]
	
	\subsection*{5. Applications in control engineering}
	\begin{itemize}
		\item Solving differential equations: convert to algebra, solve, then invert.
		\item Transfer functions: $G(s) = Y(s)/U(s)$.
		\item Stability analysis: study poles of $F(s)$ or $G(s)$.
		\item Frequency response: substitute $s = j\omega$ to obtain amplitude and phase.
	\end{itemize}
	
	\subsection*{6. Interpretation}
	The Laplace transform is used because it converts differential equations into algebraic equations, while also providing a framework to analyze how signals of different shapes (especially exponentials and sinusoids) propagate through systems.
	
	\bigskip
	
	\subsection*{Example A: RC Circuit}
	For a series RC circuit with input $u(t)$ and output $y(t) = v_C(t)$ across the capacitor:
	$$
	RC \frac{dy(t)}{dt} + y(t) = u(t).
	$$
	Taking the Laplace transform with zero initial condition:
	$$
	(RCs + 1) Y(s) = U(s).
	$$
	The transfer function is
	$$
	G(s) = \frac{Y(s)}{U(s)} = \frac{1}{RCs+1}.
	$$
	
	\subsection*{Example B: Mass–Spring–Damper}
	Equation of motion:
	$$
	M \frac{d^2 x(t)}{dt^2} + B \frac{dx(t)}{dt} + Kx(t) = F(t),
	$$
	where $x(t)$ is displacement and $F(t)$ is force input.  
	Laplace transform with zero initial conditions:
	$$
	(Ms^2 + Bs + K) X(s) = F(s).
	$$
	The transfer function is
	$$
	G(s) = \frac{X(s)}{F(s)} = \frac{1}{Ms^2 + Bs + K}.
	$$
	
	\subsection*{7. Electrical–Mechanical Analogy}
	\[
	\begin{array}{|c|c|}
		\hline
		\text{Electrical quantity} & \text{Mechanical analogue} \\
		\hline
		\text{Voltage } u(t) & \text{Force } F(t) \\
		\text{Current } i(t) & \text{Velocity } \dot{x}(t) \\
		\text{Capacitor } C & \text{Spring constant } K \\
		\text{Resistor } R & \text{Viscous damper } B \\
		\text{Inductor } L & \text{Mass } M \\
		\hline
	\end{array}
	\]
	
	
			
		\section*{Signals, Impulses, and Convolution: A Unified Note}
		
		\subsection*{1. The Unit Step}
		
		\textbf{Continuous-time:}
		$$
		u(t) =
		\begin{cases}
			1, & t \geq 0, \\
			0, & t < 0.
		\end{cases}
		$$
		
		\textbf{Discrete-time:}
		$$
		u[n] =
		\begin{cases}
			1, & n \geq 0, \\
			0, & n < 0.
		\end{cases}
		$$
		
		The step acts as a ``switch'' that turns signals on at $t=0$ or $n=0$.  
		It also delimits integrals/sums, enforcing \emph{causality}.
		
			
		\subsection*{2. The Unit Impulse}
		
		\textbf{Continuous-time (Dirac delta):}
		$$
		\delta(t) =
		\begin{cases}
			0, & t \neq 0, \\
			\text{undefined at } t=0, \;\; \int_{-\infty}^\infty \delta(t)\,dt = 1.
		\end{cases}
		$$
		It is a \emph{distribution}, defined by the sifting property
		$$
		\int_{-\infty}^{\infty} f(t)\,\delta(t)\,dt = f(0).
		$$
		
		\textbf{Discrete-time (Kronecker delta):}
		$$
		\delta[n] =
		\begin{cases}
			1, & n=0, \\
			0, & n \neq 0,
		\end{cases}
		$$
		with the sifting property
		$$
		\sum_{k=-\infty}^\infty f[k]\delta[k] = f[0].
		$$
		
			
		\subsection*{3. Relation Between Step and Impulse}
		
		\textbf{Continuous-time:}
		$$
		\frac{d}{dt}u(t) = \delta(t), 
		\qquad
		\int_{-\infty}^{t} \delta(\tau)\, d\tau = u(t).
		$$
		
		\textbf{Discrete-time:}
		$$
		\nabla u[n] = u[n] - u[n-1] = \delta[n],
		\qquad
		\sum_{k=-\infty}^{n} \delta[k] = u[n].
		$$
		
		Thus, the impulse marks the ``jump'' in the step, and the step accumulates impulses.
		
			
		\subsection*{4. Example: Finite Step and First Difference}
		
		Let
		$$
		u[n] =
		\begin{cases}
			1, & 0 \leq n \leq 3, \\
			0, & \text{otherwise}.
		\end{cases}
		$$
		
		Then
		\begin{align*}
			\nabla u[0] &= u[0]-u[-1] = 1-0 = 1, \\
			\nabla u[1] &= u[1]-u[0] = 1-1 = 0, \\
			\nabla u[2] &= 1-1=0, \\
			\nabla u[3] &= 1-1=0, \\
			\nabla u[4] &= 0-1=-1.
		\end{align*}
		
		So
		$$
		\nabla u[n] = \{1,0,0,0,-1\},
		$$
		marking where the sequence steps up and down.
		
			
		\subsection*{5. Convolution}
		
		\textbf{Continuous-time:}
		$$
		(f*g)(t) = \int_{-\infty}^{\infty} f(\tau)\,g(t-\tau)\,d\tau.
		$$
		
		\textbf{Discrete-time:}
		$$
		(f*g)[n] = \sum_{k=-\infty}^{\infty} f[k]\, g[n-k].
		$$
		
		Interpretation:
		\begin{itemize}
			\item $g(t-\tau)$ (or $g[n-k]$) is a time-reversed and shifted copy of $g$.
			\item The integral (or sum) measures the overlap between $f$ and the shifted copy of $g$.
			\item Each output is the accumulated contribution of past inputs, weighted by $g$.
		\end{itemize}
		
			
		\subsection*{6. Properties of Convolution}
		
		\begin{itemize}
			\item \textbf{Commutativity:}
			$$
			(f*g)(t) = \int f(\tau)g(t-\tau)\,d\tau 
			= \int g(\tau)f(t-\tau)\,d\tau 
			= (g*f)(t).
			$$
			
			\item \textbf{Causality:}  
			If $f(t)$ and $g(t)$ are causal (multiplied by $u(t)$), then $(f*g)(t)$ is also causal, since the overlap begins at $t=0$.
			
			\item \textbf{Linearity:}  
			$$
			f*(ag+bh) = a(f*g) + b(f*h).
			$$
		\end{itemize}
		
		
		\subsection*{7. Convolution Theorem (Laplace Domain)}
		
		Let
		$$
		F(s) = \int_0^\infty f(t)e^{-st}\,dt,
		\qquad
		G(s) = \int_0^\infty g(t)e^{-st}\,dt.
		$$
		
		Then
		$$
		\mathcal{L}\{(f*g)(t)\} = F(s)G(s).
		$$
		
		Thus:
		\begin{itemize}
			\item Convolution in time $\;\longleftrightarrow\;$ Multiplication in $s$-domain.
			\item Multiplication in time $\;\longleftrightarrow\;$ Convolution in $s$-domain.
		\end{itemize}
		
		This simplifies analysis: a complicated time-domain accumulation becomes a simple algebraic product.
		
		
		\subsection*{8. Mental Map}
		
		\begin{itemize}
			\item The unit step $u(t)$ switches signals on and enforces causality.
			\item The unit impulse $\delta(t)$ (or $\delta[n]$) is the ``atom'' of signals:  
			- It marks changes (as derivative/difference of $u$).  
			- It builds signals (superposition of shifted impulses).
			\item Convolution is the \emph{memory mechanism}: the output is the accumulation of past inputs, each leaving an echo shaped by the impulse response.
			\item In the Laplace domain, this accumulation appears as simple multiplication, turning system analysis into algebra.
		\end{itemize}
		
	
	
	
\end{document}
